%% quantmsdiann manuscript — Molecular & Cellular Proteomics (MCP)
%% Article type: Software
%%
%% Build: cd paper && make pdf
%% Figures: MCP column widths via mcpfigures.sty; SVG sources in ../analysis/figures/
%%
%% TODO before submission — author-fillable placeholders:
%%   declarations.tex: Competing interests, Funding, Author contributions, Acknowledgements (currently draft placeholders)
%%
%% Pipeline version: v2.2.0 (release codename "Chongqing", 2026-06-16) — filled in methods and declarations

\documentclass[preprint,12pt]{elsarticle}
%% For journal submission, switch to double-column layout, e.g.:
%% \documentclass[final,3p,times,twocolumn]{elsarticle}

\usepackage[utf8]{inputenc}
\usepackage[T1]{fontenc}
\usepackage{microtype}
\usepackage{grffile}
\usepackage{amsmath}
\usepackage{enumitem}
\usepackage{url}
\usepackage{hyperref}
\usepackage[margin=2cm]{geometry}
\usepackage{mcpfigures}
\usepackage{orcidlink}
\input{macros}

%% Keep long URLs, container names, paths and command tokens inside the
%% margins: allow URLs to break anywhere, and give the line breaker extra
%% slack so no \texttt{} token overruns the text block.
\usepackage{xurl}
\Urlmuskip=0mu plus 1mu
\emergencystretch=3em

\journal{Molecular \& Cellular Proteomics}

\begin{document}

\begin{frontmatter}

\title{quantmsdiann: a scalable SDRF-driven DIA-NN workflow for reanalysis of single-cell, spatial, and bulk proteomics datasets}

\author[aff1]{Qi-Xuan Yue\,\orcidlink{0009-0004-8120-9339}}
\author[aff1]{Yufei Shen}
\author[aff3,aff4]{Asier Larrea-Sebal\,\orcidlink{0000-0001-9107-4299}}
\author[aff2]{Chengxin Dai\,\orcidlink{0000-0001-6943-5211}}
\author[aff5]{Henry Webel\,\orcidlink{0000-0001-8833-7617}}
\author[aff7,aff12,aff13]{Jose Nimo\,\orcidlink{0000-0002-1565-7799}}
\author[aff7,aff13]{Fabian Coscia\,\orcidlink{0000-0002-2244-5081}}
\author[aff8]{Orhun Kok\,\orcidlink{0009-0002-9764-9094}}
\author[aff8,aff9]{Nikolai Slavov\,\orcidlink{0000-0003-2035-1820}}
\author[aff3]{Juan Antonio Vizca\'{\i}no\,\orcidlink{0000-0002-3905-4335}}
\author[aff10,aff11]{Timo Sachsenberg\,\orcidlink{0000-0002-2833-6070}}
\author[aff1]{Mingze Bai\,\orcidlink{0000-0002-9782-2056}}
\author[aff6]{Vadim Demichev\,\orcidlink{0000-0002-2424-9412}}
\author[aff3]{Yasset Perez-Riverol\,\orcidlink{0000-0001-6579-6941}\corref{cor1}}

\ead[cor1]{yasset@ebi.ac.uk}
\cortext[cor1]{Corresponding author.}

\affiliation[aff1]{organization={Chongqing Key Laboratory of Big Data for Bio Intelligence, Chongqing University of Posts and Telecommunications}, city={Chongqing}, country={China}}
\affiliation[aff2]{organization={State Key Laboratory of Proteomics, Beijing Proteome Research Center, National Center for Protein Sciences (Beijing), Beijing Institute of Life Omics}, city={Beijing}, country={China}}
\affiliation[aff3]{organization={European Molecular Biology Laboratory, European Bioinformatics Institute (EMBL-EBI), Wellcome Genome Campus}, city={Hinxton, Cambridge}, country={UK}}
\affiliation[aff4]{organization={Department of Biochemistry and Molecular Biology, Universidad del Pa\'{i}s Vasco UPV/EHU}, city={Bilbao}, postcode={48080}, country={Spain}}
\affiliation[aff5]{organization={Novo Nordisk Foundation Center for Protein Research, Faculty of Health and Medical Sciences, University of Copenhagen}, city={Copenhagen}, country={Denmark}}
\affiliation[aff6]{organization={Quantitative Proteomics laboratory, Charit\'{e} -- Universit\"{a}tsmedizin Berlin}, city={Berlin}, country={Germany}}
\affiliation[aff7]{organization={Spatial Proteomics Group, Max-Delbr\"{u}ck-Center for Molecular Medicine in the Helmholtz Association (MDC)}, city={Berlin}, country={Germany}}
\affiliation[aff8]{organization={Department of Bioengineering, Department of Biology, Department of Chemistry and Chemical Biology, Single Cell Proteomics Center, and Barnett Institute for Chemical and Biological Analysis, Northeastern University}, city={Boston, Massachusetts}, country={USA}}
\affiliation[aff9]{organization={Parallel Squared Technology Institute}, city={Watertown, Massachusetts}, country={USA}}
\affiliation[aff10]{organization={Department of Computer Science, Applied Bioinformatics, University of T\"{u}bingen}, city={T\"{u}bingen}, country={Germany}}
\affiliation[aff11]{organization={Institute for Bioinformatics and Medical Informatics, University of T\"{u}bingen}, city={T\"{u}bingen}, country={Germany}}
\affiliation[aff12]{organization={Charit\'{e} – Universit\'{a}tsmedizin Berlin, corporate member of Freie Universit\'{a}t Berlin and Humboldt-Universit\'{a}t zu Berlin},city={Berlin}, country={Germany}}
\affiliation[aff13]{organization={German Cancer Consortium (DKTK), Partner Site Berlin}, city={Heidelberg}, country={Germany}}


\begin{abstract}
Public proteomics archives now hold thousands of data-independent acquisition (DIA) datasets, but reusing them is difficult: each was processed with a different software configuration, and most lack standardised metadata. We present \quantmsdiann{}, an open-source Nextflow/nf-core workflow that runs \texttt{DIA-NN} in parallel across cloud and high-performance computing (HPC) infrastructure, guided by the experimental design declared in SDRF format. The workflow ships container-pinned profiles for each supported \texttt{DIA-NN} version through Docker and Singularity, needs no manual installation or dependency resolution, and reads all major vendor formats and the HUPO-PSI mzML standard, natively where the \texttt{DIA-NN} release supports it and by automatic conversion otherwise. It exports harmonised quantification tables as MSstats input, the Quantitative Proteomics eXchange format (\qpx{}), and a \texttt{pmultiqc} quality-control report. The parallel design reanalyses a 2{,}300-run single-cell dataset in 2.2~hours on 300 HPC nodes. On the two ProteoBench modules with matched community submissions, distributed \quantmsdiann{} reproduces the single-machine quantification, and identification depth increases with each \texttt{DIA-NN} release across all four modalities. Reanalysing public DIA deposits with a current \texttt{DIA-NN} build matches or exceeds the originally deposited counts, recovering up to 59\% more protein groups, with the largest gains on deposits processed with older or non-\texttt{DIA-NN} engines and smaller gains where a recent \texttt{DIA-NN} release was already used. \quantmsdiann{} is a step towards scalable, reproducible reanalysis of the growing DIA data archive and a foundation for large-scale DIA meta-analysis and atlas building. It is available at \url{https://github.com/bigbio/quantmsdiann}.
\end{abstract}

\begin{keyword}
DIA-NN \sep single-cell proteomics \sep spatial proteomics \sep phosphoproteomics \sep plexDIA \sep data-independent acquisition \sep Nextflow \sep nf-core \sep SDRF \sep reproducibility \sep QPX \sep pmultiqc
\end{keyword}

\end{frontmatter}

\section{Introduction}
\label{sec:introduction}
%% [begin sections/introduction.tex]
Data-independent acquisition (DIA) is a rapidly growing mass-spectrometry strategy for proteomics at scale, now approaching parity with data-dependent acquisition among public proteomics deposits, and \texttt{DIA-NN}~\cite{demichev2020} is among the most widely adopted analysis engines across timsTOF~\cite{meier2020,brunner2022}, Astral~\cite{bubis2025astral}, and Orbitrap platforms. Recent benchmarks confirm \texttt{DIA-NN}'s competitiveness for single-cell DIA workflows~\cite{wang2025}, and recent releases have expanded support for plexDIA~\cite{derks2023, derks2023strategies}, timsTOF parallel accumulation-serial fragmentation (PASEF), and vendor-native raw-file reading. At the same time, DIA dataset sizes and run counts are growing rapidly~\cite{perezriverol2025pride}, and the proteomics community is increasingly interested in reanalysing public deposits to recover additional identifications, integrate across cohorts, and build large-scale DIA atlases~\cite{dai2024quantms,lautenbacher2022proteomicsdb}.
\texttt{DIA-NN}'s command-line interface, simplicity, and analytical power make it a natural engine for mining public archives, but realising it at scale requires distributing the computation: extracting more information from public data means reprocessing cohorts of hundreds to thousands of runs fast enough that reanalysis can be iterative rather than a one-off effort. The same need extends beyond reanalysis: laboratories generating their own large cohorts increasingly have access to cloud and high-performance computing (HPC) infrastructure and benefit from distributing the analysis rather than running it on a single workstation. The steps of a \texttt{DIA-NN} run are heterogeneous and map naturally onto heterogeneous hardware: many lightweight per-file passes that fit small nodes, and a few memory- and compute-heavy cohort-level steps that need large nodes. Efficient large-scale analysis therefore depends on software containers and workflow engines that can place each task on appropriate cloud or HPC resources~\cite{perezriverol2020scalable}. Despite this analytical maturity, most \texttt{DIA-NN} users run the tool as a desktop application, which limits its scalability for en masse processing of very large datasets.

In 2024 we published bigbio/quantms~\cite{dai2024quantms}, a Nextflow pipeline that standardised multi-engine DIA and DDA analyses behind a Sample and Data Relationship Format (SDRF)~\cite{dai2021sdrf} input on the nf-core framework~\cite{ewels2020nfcore}. quantms was, however, built around \texttt{DIA-NN}~1.8.1 with a fixed parameter set, and its DDA/DIA combination within the same workflow made extending DIA-specific functionality difficult. New \texttt{DIA-NN} parameters conflicted with DDA options, versions were hard-coded, sample-level acquisition parameters were not properly propagated to per-run \texttt{DIA-NN} invocations, and raw-file conversion was mandatory. Iterative reanalysis at archive scale requires a workflow that combines the reproducibility of nf-core, the sensitivity of current \texttt{DIA-NN} releases, the metadata discipline of SDRF, and a parallelisation strategy capable of processing thousands of raw files.

Here, we present \quantmsdiann{}, an independent SDRF-driven Nextflow pipeline dedicated to \texttt{DIA-NN} that extends DIA-specific features beyond what the original bigbio/quantms pipeline could accommodate. While \quantmsdiann{} shares its SDRF-first design philosophy and parts of the configuration-generation tooling with quantms, its workflow, container set, and module library are developed and released independently, with container-pinned profiles for every \texttt{DIA-NN} version selectable at runtime under Docker or Singularity. Downstream, it exports a harmonised \qpx{} (Quantitative Proteomics eXchange) Parquet archive alongside the standard \texttt{DIA-NN} reports and MSstats input. We demonstrate the pipeline along three axes: ProteoBench~\cite{devreese2025proteobench} benchmarks across four DIA modalities and successive \texttt{DIA-NN} releases; independent reanalyses of public DIA deposits spanning bulk cell lines (NCI-60, ProCan), single-cell datasets, a spatial Deep Visual Proteomics cohort, and a phosphoproteomics cohort; and a pan-cohort of five public DIA cell-line reanalyses integrated under a single SDRF-driven invocation. Across the single-cell and bulk cell-line cohorts, reanalysis with the current \texttt{DIA-NN} build recovers up to 59\% more protein groups, with the largest gains on the oldest deposits processed with older or non-\texttt{DIA-NN} engines and smaller gains where a recent \texttt{DIA-NN} release was already used, so reprocessing archived data with an up-to-date engine yields measurably more from each deposit.
%% [end introduction]

\section{Experimental Procedures}
\label{sec:methods}
%% [begin sections/methods.tex]
\subsection{Pipeline implementation}
\label{sec:impl}

\quantmsdiann{} is implemented in Nextflow DSL2 ($\geq$25.10.4) and follows nf-core community standards~\cite{ewels2020nfcore} for pipeline structure, testing, documentation, and continuous integration. It supports both \texttt{DIA-NN} library-free analysis (predicting the spectral library in-silico from a FASTA reference) and library-based analysis against a user-supplied spectral library; all analyses in this study used the library-free mode; the reanalysis cohorts used the UniProt human reference proteome~\cite{uniprot2025} (UP000005640, downloaded March~2026), while the ProteoBench HYE benchmarks used the corresponding three-species (human, yeast, \textit{E.\ coli}) FASTA. Each computational step is packaged as a versioned container image (Docker, Singularity, or Apptainer). Open-source community tools (ThermoRawFileParser~\cite{hulstaert2020}, pridepy~\cite{kamatchinathan2025pridepy}, \qpx{}, and the \texttt{wiffconverter} SCIEX reader) are pulled from public registries (BioContainers~\cite{biocontainers2017} or \texttt{ghcr.io/bigbio}).
Because the \texttt{DIA-NN} license restricts redistribution, only version~1.8.1 is shipped as an image (the default, from BioContainers); releases from 1.9 onward (1.9.2--2.5.1, including the 2.5.1-enterprise build) are built locally from the per-release recipes in the companion \texttt{quantms-containers} repository (\url{https://github.com/bigbio/quantms-containers}) and tagged so the workflow selects them automatically (Supplementary Note~1). The repository at \url{https://github.com/bigbio/quantmsdiann} is organised as a standard nf-core pipeline with main workflow (\texttt{main.nf}), subworkflows, modules, test data and CI configurations, and documentation. Pipeline version v2.2.0 (release codename ``Chongqing'', 2026-06-16) was used for the analyses reported in this paper.

\subsection{SDRF extensions for quantmsdiann}
\quantmsdiann{} accepts SDRF~\cite{dai2021sdrf} as primary input.
SDRF parsing is performed using the sdrf-pipelines Python package developed within the quantms project, which validates the SDRF against the proteomics-specific subset of EFO, PSI-MS, and PRIDE ontology terms and emits a normalised per-sample channel for downstream Nextflow processes. To express DIA- and \texttt{DIA-NN}-specific acquisition settings at the sample level, we extended the SDRF-Proteomics format (version~1.1.0~\cite{dai2021sdrf}; specification and templates at \url{https://github.com/bigbio/proteomics-sample-metadata}) with additional comment columns: the precursor (MS1) and fragment (MS2) $m/z$ scan ranges and, for multiplexed experiments, the plexDIA channel definitions and their isotopologue labels. Because these settings are declared per sample, independent datasets acquired with different scan-range schemes are each parametrised from their own metadata and can be processed together in one invocation, with no change to the workflow. We added a \texttt{convert-diann} command to sdrf-pipelines that reads these extended columns together with the standard fields (enzyme, fixed and variable modifications, instrument, precursor mass tolerance, fragment mass tolerance, and missed-cleavage settings) and writes a per-run \texttt{DIA-NN} command-line configuration; where SDRF columns are absent, \texttt{DIA-NN} defaults appropriate for the inferred instrument are applied.

\subsection{Raw file handling and \texttt{DIA-NN} execution}
A key advance over its predecessor quantms is that native vendor (raw) formats are first-class inputs in \quantmsdiann{}: whereas quantms converted every raw file to mzML up front, \quantmsdiann{} reads each instrument's native format directly and converts only when a particular format or \texttt{DIA-NN} version requires it. Thermo \texttt{.raw} or timsTOF (.d or .d.zip) files are read natively by \texttt{DIA-NN}~$\geq$~2.1.0; for DIA-NN versions lower than 2.1.0 conversion to mzML is performed with ThermoRawFileParser~\cite{hulstaert2020} (\texttt{--mzml\_convert}); Bruker \texttt{.d} folders are ingested natively, with compressed forms (\texttt{.d.tar}/\texttt{.zip}/\texttt{.gz}) decompressed first; and Sciex \texttt{.wiff}~/~\texttt{.wiff2} are converted to mzML with the open-source \texttt{wiffconverter} (\url{https://github.com/bigbio/wiffconverter}). Generic \texttt{.dia} and pre-converted \texttt{.mzML} inputs pass through unchanged, with optional OpenMS re-indexing (\texttt{--reindex\_mzml}).

\subsection{Quality control and \qpx{} export}
Per-run and per-cohort quality control reports are generated by \texttt{pmultiqc}~\cite{yue2026pmultiqc}, a proteomics extension of MultiQC~\cite{ewels2016multiqc} that summarises \texttt{DIA-NN}-specific metrics including peptide and protein counts, missed cleavages, charge-state distributions, precursor and fragment mass accuracy, and per-run identification yield. The \qpx{} archive is produced by the \texttt{qpxc} CLI from the \texttt{qpx} package (\url{https://github.com/bigbio/qpx}), which writes Parquet files for identifications (peptide-spectrum matches (PSMs), peptides, protein groups, features), metadata (sample, run, ontology), and provenance, and optionally exports MuData (\texttt{.h5mu}) containers for scverse-based~\cite{virshup2023scverse} downstream analysis. The aim of exporting to \qpx{} is to bundle the SDRF, identifications, quantifications, and pipeline provenance into a single queryable artefact for reanalysis and long-term archival, and to provide a standardised input for downstream tools that support the format (e.g.\ MSstats, Perseus, or custom scripts).

\subsection{Datasets benchmarked and compute environment}
\label{sec:datasets}
Supplementary Note~2 (Table~S1) lists all benchmark and reanalysis datasets with their instrument, acquisition mode, and deposition year, and an SDRF prepared from the original sample metadata. They fall into four experiment classes: \emph{ProteoBench} DIA benchmarks (modules~5, 7, 9, and 10); \emph{single-cell and spatial proteomics}, comprising the HeLa single-cell cohort PXD046357 (Orbitrap Astral), the oocyte plexDIA cohort MSV000093870, the PXD064049 spatial Deep Visual Proteomics cohort, and the PXD071075 single-cell scalability sweep; \emph{bulk cell-line panels} (PXD003539, PXD030304, PXD004701, PXD041421, and PXD017199); and \emph{phosphoproteomics} (PXD034128, PXD034623, and PXD049692).
For each cell-line reanalysis cohort, the original published quantification was downloaded from a public repository as the comparator. \quantmsdiann{} counts follow a single global rule (protein groups at \texttt{Lib.PG.Q.Value}~$\leq 0.01$, precursors at \texttt{Lib.Q.Value}~$\leq 0.01$, decoys dropped, no contaminant or positive-quantity filter); the like-for-like-threshold rule applied to each original headline count is described in Supplementary Note~3. Runtime benchmarks were run on the EMBL-EBI HPC cluster (LSF scheduler), with per-task CPU and memory dispatched by Nextflow and memory profiled from the Nextflow trace report (\texttt{-with-trace}).

\subsection{Benchmarking and cross-study comparison}
\label{sec:benchmarking}
A distributed workflow is only useful if it reproduces a single-machine \texttt{DIA-NN} run, so we benchmarked \quantmsdiann{} against the ProteoBench community submissions~\cite{devreese2025proteobench} on the four public DIA modules (Orbitrap Astral, single-cell Astral, timsTOF diaPASEF, and SCIEX ZenoTOF~7600). All \quantmsdiann{} runs were library-free (the library predicted in-silico from the FASTA), so for a like-for-like comparison we kept only predicted-from-FASTA \texttt{DIA-NN} community submissions and report the oldest supported release (1.8.1) and the current 2.5.1-enterprise build at \texttt{DIA-NN}'s recommended per-version q-values under a uniform 1\% global precursor cut-off. For the public reanalyses, each deposit's original published quantification served as the comparator, counted under that study's own threshold.

Cross-engine identification counts carry three caveats that apply to every comparison in the Results: differences in FDR calibration across search-engine generations, the use of each original study's own counting criterion, and possible differences in protein sequence database and search settings between the deposited analysis and the reanalysis. These are set out where the comparisons are presented (Results, ``The value of reanalysing public DIA deposits at scale'').
%% [end methods]

\section{Results}
\label{sec:results}
%% [begin sections/results.tex]
\subsection{\quantmsdiann{} pipeline architecture}
\label{sec:architecture}

Built on the nf-core framework~\cite{ewels2020nfcore}, \quantmsdiann{} accepts a single SDRF file as primary input, which encodes per-sample metadata (cell line or tissue origin, factor values, instrument, post-translational modifications) and per-run file references. We extended the SDRF with DIA-relevant columns (MS1 and MS2 scan ranges, and plexDIA channel definitions and labels) and added an adapter to sdrf-pipelines~\cite{dai2021sdrf} (\texttt{convert-diann}) that translates each run's declared metadata into a per-run \texttt{DIA-NN} command-line configuration. Every run is therefore parametrised from its own metadata with no manual configuration or parameter-harmonisation step, removing a common source of analysis-script divergence between groups: independent datasets with different scan-range schemes, or a multiplexed plexDIA design, are each analysed directly from their declared SDRF settings without any workflow change.
The workflow comprises four mandatory and three optional modules organised around the \texttt{DIA-NN} analysis core (Fig.~\ref{fig:architecture}), alternating between per-file parallel stages (red; raw-file preparation, preliminary calibration, and individual search) and collective stages (green; in-silico library generation, empirical library assembly, final quantification, MSstats conversion, and \texttt{pmultiqc}~\cite{yue2026pmultiqc} quality-control reporting) operating on the full cohort.

\quantmsdiann{} follows \texttt{DIA-NN}'s distributed multi-pass flow, alternating parallel per-file work with serial cohort-wide steps (Fig.~\ref{fig:architecture}). Each raw file is first prepared \emph{in parallel} (vendor-format conversion, decompression, or indexing) where needed; conversion to mzML is mandatory for early \texttt{DIA-NN} versions and for some vendor formats such as SCIEX \texttt{.wiff}. A spectral library is then predicted in-silico from the FASTA once (\emph{serial}; skipped when a library is supplied), and every file is calibrated against it in a \emph{parallel} first pass. The first-pass results are merged into a single empirical library (\emph{serial}); each file is searched again against that library \emph{in parallel}; and a final consolidation pass merges all per-file results into the cohort quantification, applying match-between-runs, normalisation, and protein inference (\emph{serial}). MSstats conversion, \texttt{pmultiqc} reporting, and \qpx{} export then run once at the end (\emph{serial}). The two per-file search passes carry essentially all of the compute and scale across cluster nodes, driving the throughput reported below; the serial steps set a largely fixed wall-clock floor.
Outputs include the standard \texttt{DIA-NN} report tables: \texttt{report.tsv} (or \texttt{report.parquet}), \texttt{report.pr\_matrix.tsv}, and \texttt{report.pg\_matrix.tsv}, alongside MSstats-formatted input and a \texttt{pmultiqc}~\cite{yue2026pmultiqc} quality-control report. A harmonised \qpx{} archive bundles identifications, metadata, and provenance into a Parquet-based artefact, queried directly with DuckDB and, for the quantification layers, exported to MuData (\texttt{.h5mu}) for scverse-compatible~\cite{virshup2023scverse} downstream analysis.

Two design choices make \quantmsdiann{} a reanalysis platform rather than a single-analysis tool. Because every \texttt{DIA-NN} release ships as a container-pinned profile selectable at runtime, the same SDRF can be analysed with any supported version (1.8.1 through 2.5.1, including the enterprise build) without reinstallation. Because the per-file work is parallelised, a single invocation scales from six-run benchmarks to cohorts of thousands of runs. We next establish that this distributed execution matches a single-machine \texttt{DIA-NN} run, then use it to reanalyse public deposits with current releases.

\begin{figure}[htbp]
  \centering
  \mcpincludegraphics[width=0.82\linewidth,keepaspectratio]{manuscript/workflow_panel}
  \caption[The quantmsdiann workflow]{%
    \textbf{The \quantmsdiann{} workflow.} Subway diagram: mandatory modules as
    filled circles, optional/skippable modules as open circles. Per-file parallel
    stages (red; raw-file preparation, preliminary calibration, individual search)
    run as one task per raw file, whereas collective stages (green; in-silico
    library generation, empirical library assembly, final quantification, MSstats
    conversion, \texttt{pmultiqc} reporting, and \qpx{} export) run once over the
    full cohort.
  }
  \label{fig:architecture}
\end{figure}

\subsection{quantmsdiann throughput, scaling, and analytical equivalence}
\label{sec:throughput}

To characterise how \quantmsdiann{} scales, we reanalysed an entire single-cell experiment of \textasciitilde 2{,}300 raw files (PXD071075, Orbitrap Eclipse) across cluster sizes from 10 to 300 nodes and measured end-to-end wall-clock time (Fig.~\ref{fig:validation}a). Because the per-file work is parallelised across nodes, a single SDRF-driven invocation completed the whole cohort within hours, falling from 37.4~h at 10 nodes to 2.2~h at 300. Scaling stays close to the ideal $1{/}N$ reference up to \textasciitilde 100 nodes and then bends away, with a practical knee near 200 nodes beyond which added nodes return progressively less. This floor reflects the serial cohort-level steps that run once regardless of node count and contention for shared-cluster I/O and scheduler slots, so part of it is specific to this allocation rather than intrinsic to the workflow. The practical guidance is to size the cluster to the desired turn-around time rather than to maximum throughput.

We then compared time-to-finish across every dataset in the study (Fig.~\ref{fig:validation}b), each reported net of the one-time in-silico library-prediction step, which cannot be parallelised and whose cost depends on the library-generation parameters. (The sole exception is PXD003539, for which the Nextflow execution trace was incomplete and the library step could not be separated out, so its bar conservatively includes it, a $\sim$0.4~h overstatement.) Across cohorts spanning three orders of magnitude in file count, wall-clock stayed within hours because the per-file passes absorb additional files in parallel: the 5{,}798-run ProCan cohort (PXD030304) finished in 6.2~h, in the same band as cohorts orders of magnitude smaller. Time-to-finish is therefore governed by per-run search cost and cohort-level consolidation rather than by file count; the longest reanalysis was a 206-run cohort on an older Q~Exactive instrument (PXD017199, 11.4~h), exceeding the far larger ProCan cohort. Peak memory during consolidation stayed within the 64--128~GB available on commodity HPC nodes, so no specialist high-memory hardware is required. The plexDIA and phosphoproteomics cohorts, processed to demonstrate workflow support rather than for cross-engine benchmarking, likewise complete within hours (Fig.~\ref{fig:validation}b), with per-step runtime and resource envelopes for the benchmark and sweep cohorts in Supplementary Note~4 (Figs.~S1--S2). To confirm the workflow is not tied to one infrastructure, an independent group ran the workflow on a separate (non-EBI) SLURM cluster with the same container-pinned profiles and observed the same per-step runtime structure (Supplementary Fig.~S1b).

With scaling characterised, we asked whether distributed, SDRF-driven execution matches a single-machine \texttt{DIA-NN} run. On the two ProteoBench modules with predicted-from-FASTA \texttt{DIA-NN} community submissions (Orbitrap Astral, $n=15$; timsTOF diaPASEF, $n=8$), \quantmsdiann{} reproduces the expected HYE per-species log$_2$ fold-changes, and across both releases its measured ratios sit within the spread of the standalone desktop submissions (Fig.~\ref{fig:validation}c). Distributed \quantmsdiann{} thus matches desktop \texttt{DIA-NN} quantification accuracy on these two modules, within the dispersion of the community submissions.

Identification depth, by contrast, increases with each \texttt{DIA-NN} version on every module. Per-run protein groups rise consistently from 1.8.1 through 2.5.1 to the enterprise build across all four DIA modalities (Fig.~\ref{fig:validation}d), a gain of 10--12\%: for example, from 7{,}305 to 8{,}190 protein groups per run on the single-cell Astral module (Module~9) and from 10{,}192 to 11{,}432 on Orbitrap Astral (Module~7). Precursors and the number of proteins quantified in every run increase in parallel (per-version precursor counts in Supplementary Note~5; the complete-profile protein-group breakdown in Fig.~S3). This tracks the broader ProteoBench community trend~\cite{devreese2025proteobench} and motivates reanalysis of public deposits with the current build.

\begin{figure}[htbp]
  \centering
  \mcpincludegraphics[width=\textwidth,keepaspectratio]{manuscript/fig2_validation}
  \caption[Scaling and ProteoBench accuracy]{%
    \textbf{Runtime scaling and ProteoBench quantification accuracy.}
    (\textbf{a})~Workflow wall-clock versus number of cluster nodes (Nextflow
    \texttt{executor.queueSize}) on the PXD071075 single-cell cohort (log--log;
    one run per point). The dashed grey line is ideal $1{/}N$ scaling anchored at
    10 nodes; parallel efficiency is printed under each dot. The measured curve
    rises above the ideal line beyond \textasciitilde 100 nodes (parallel
    efficiency falling from 100\% to 58\%), reflecting fixed serial cohort-level
    steps and shared-cluster contention on this EMBL-EBI allocation.
    (\textbf{b})~Wall-clock per reanalysis, one bar per dataset, ordered by cohort
    size and coloured by instrument family; the dataset-type metadata is shown by
    the colour of each cohort's y-axis label (descriptive only: bar height reflects
    cohort size and parallelisation, not sample type per se).
    (\textbf{c})~ProteoBench accuracy per HYE species: observed log$_2$ fold-change
    against the expected ratio (dashed) for the two DIA modules with
    predicted-from-FASTA \texttt{DIA-NN} submissions (Orbitrap Astral, Module~7;
    timsTOF diaPASEF, Module~5, PXD062685). Standalone community runs (all versions) form
    the grey strip and box; \quantmsdiann{} points are coloured by instrument and
    shaped by \texttt{DIA-NN} version (circle, 1.8.1; triangle, 2.5.1-enterprise),
    tracking the community distribution, i.e.\ distributed orchestration does not
    change quantification accuracy.
    (\textbf{d})~Protein-group identification depth on ProteoBench across three
    \texttt{DIA-NN} configurations (1.8.1, 2.5.1, and the 2.5.1-enterprise build)
    and four DIA modalities, as protein groups per run (averaged over the six
    replicates) at a 1\% run-specific protein-group q-value. Depth increases with
    each release on every module, tracking the accuracy trend in (c); the
    complete-profile companion (protein groups quantified in every run) and the
    per-version precursor breakdown are in Supplementary Note~5, Fig.~S3.
  }
  \label{fig:validation}
\end{figure}

\subsection{Single-cell DIA reanalysis improves with current \texttt{DIA-NN} releases}
\label{sec:single-cell}

Single-cell proteomics has rapidly advanced, with significant contributions from computational methods~\cite{slavov2025Review} and community recommendations for performing, benchmarking, and reporting experiments~\cite{White_Paper}. To assess \quantmsdiann{} on low-input data, where sensitivity and match-between-runs matter most, we reanalysed two label-free Orbitrap Astral single-cell cohorts, HeLa (PXD046357) and a larger A549/H460 cohort of 146 single cells (PXD049412), through the workflow under \texttt{DIA-NN}~1.8.1 and the current 2.5.1-enterprise build, both parametrised from the same SDRF (Fig.~\ref{fig:single-cell}).
Upgrading from 1.8.1 to 2.5.1-enterprise increased identifications on both cohorts (Fig.~\ref{fig:single-cell}a). On HeLa Astral, precursors and protein groups each rose by \textasciitilde 17\%, and the gain held per cell. On the deeper A549/H460 cohort the reanalysis reproduced the study's reported Astral single-cell depth~\cite{bubis2025astral}, a median of \textasciitilde 3{,}400 protein groups per cell; here the precursor gain exceeded the protein-group gain, consistent with the near-saturated protein depth of this larger cohort.
The improvement is largest in the data-completeness profile (the number of protein groups quantified across an increasing fraction of cells; Fig.~\ref{fig:single-cell}b), consistent with stronger cross-run propagation, while quantitative precision is preserved (Fig.~\ref{fig:single-cell}c). As in the cross-deposit comparisons, these version gains are identifications recovered under a current, FDR-controlled engine rather than strictly FDR-matched, and \texttt{DIA-NN} calibration in library-free, low-input mode remains an open question~\cite{wen2025fdr,krull2024singlecell}.

\begin{figure}[htbp]
  \centering
  \mcpincludegraphics[width=\linewidth,height=0.84\textheight,keepaspectratio]{manuscript/fig3_single_cell_combined}
  \caption[Single-cell DIA reanalysis]{%
    \textbf{Single-cell DIA reanalysis by \quantmsdiann{}, comparing
    \texttt{DIA-NN}~1.8.1 with the current 2.5.1-enterprise build on identical
    raw data.} Two label-free Orbitrap Astral cohorts: HeLa (PXD046357) and a
    larger A549/H460 cohort of 146 single cells (PXD049412).
    (\textbf{a})~Total identifications and per-cell depth per build: total
    precursors (left axis), total protein groups, and per-cell protein-group
    distribution (box and jittered points, right axis), with per-build percentage
    change annotated. A549/H460 reaches comparable per-cell depth on many more
    cells but is already saturated in total protein groups, so the version effect
    is concentrated in precursor depth.
    (\textbf{b})~Data-completeness curves (HeLa solid, A549/H460 dashed): protein
    groups quantified in at least $N$ cells (match-between-runs signature; right
    endpoint is the complete-profile count).
    (\textbf{c})~Per-protein coefficient of variation across cells (HeLa solid,
    A549/H460 dashed; axis truncated at CV~$=1.5$), summarising quantitative
    precision; the larger A549/H460 cohort shows the expected higher
    cell-to-cell variability.
  }
  \label{fig:single-cell}
\end{figure}

\subsection{The value of reanalysing public DIA deposits at scale}
\label{sec:reanalyses}

Public DIA deposits are typically processed once, with the search engine available at the time of deposition, and are seldom revisited, leaving recoverable identifications in raw data that has already been collected and published. Reanalysing these deposits with a current \texttt{DIA-NN} build recovers them: across every deposit reanalysed here (single-cell, bulk cell-line, and spatial), reprocessing yielded more protein groups than the originally deposited or published analysis, and more precursors wherever the original was itself \texttt{DIA-NN} (Fig.~\ref{fig:reanalysis}). This gain therefore holds across deposits rather than being specific to any single dataset.

Three caveats bound how these cross-deposit comparisons should be read. First, \emph{FDR calibration}: the comparisons assume that recent \texttt{DIA-NN} (2.5/2.5.1) exercises largely correct false-discovery-rate (FDR) control, whereas older releases (the 1.7-series and 1.8.1) and other search engines may carry substantially higher true FDR at the same nominal 1\% q-value~\cite{wen2025fdr,krull2024singlecell,frohlich2022benchmarking}, and verifying each deposit's FDR calibration is beyond the scope of this work; the reported gains should therefore be read as identifications recovered under a current, FDR-controlled engine, not as strictly FDR-matched differences. Second, \emph{counting criterion}: for the same reason we compare counts only under each original study's own threshold, and do not compare plexDIA protein groups or phosphosite numbers across engine versions. Third, \emph{database and search settings}: the reanalysis applies one harmonised FASTA and parameter set across all cohorts, whereas deposited analyses used their own databases and settings, including contaminant databases that are often not deposited (as for the spatial Deep Visual Proteomics cohort, PXD064049); residual count differences attributable to such undocumented choices cannot be fully excluded.

\begin{figure}[htbp]
  \centering
  \mcpincludegraphics[width=\textwidth,height=\mcpmaxfigheight,keepaspectratio]{manuscript/fig_reanalysis_improvement}
  \caption[Reanalysis vs original analysis]{%
    \textbf{Reanalysis with current \texttt{DIA-NN} identifies more than the
    originally deposited analyses.} Each public DIA deposit is shown as the
    originally deposited or published count (grey) versus the \quantmsdiann{}
    reanalysis (coloured by the \texttt{DIA-NN} version used: 2.5.1 or
    2.5.1-enterprise), sorted by gain.
    (\textbf{a})~Protein groups for the five reanalysed deposits, spanning
    single-cell (HeLa Astral, PXD046357), bulk cell-line (NCI-60, PXD003539;
    ProCan, PXD030304; Sun, PXD004701) and spatial Deep Visual Proteomics
    (PXD064049). The multiplexed plexDIA cohort (MSV000093870) is not shown:
    plexDIA protein-group counts are not comparable across \texttt{DIA-NN}
    versions.
    (\textbf{b})~Precursors, restricted to the two deposits with a
    \texttt{DIA-NN}-based original precursor count (HeLa Astral and spatial
    DVP); the OpenSWATH/PCT-SWATH originals provide no comparable precursor
    count.
    The \quantmsdiann{} side is counted from the per-precursor report (parquet)
    for all cohorts; the original side is the deposit's count matrix or published
    headline, except for the spatial cohort, where both the deposited
    \texttt{DIA-NN}~1.8.1 and \quantmsdiann{} counts are taken from their
    per-precursor reports under the same global rule (Supplementary Note~3), and
    over the identical 12 quantified runs on both sides (the deposited
    \texttt{DIA-NN} log loads additional library/QC runs that are not in the
    deposited quantification matrix). The
    five-cohort pan-cohort integration (per-tissue coverage) is in Supplementary
    Fig.~S6.
  }
  \label{fig:reanalysis}
\end{figure}

To show the value of running one SDRF-driven workflow across many independent deposits, we reanalysed five public cell-line DIA panels in a single \quantmsdiann{} invocation~\cite{guo2019nci60,goncalves2022procan,sun2023breast,wang2023multipro,tognetti2021breast}: NCI-60 (originally OpenSWATH), ProCan-DepMapSanger (originally \texttt{DIA-NN}~1.7.x), the Sun PCT-SWATH breast-cancer panel, the MultiPro batch-effect testbed, and Tognetti (67 breast-cancer lines); accessions are listed in the figure legend and Data Availability (Supplementary Fig.~S6).
For the two cohorts with a published DIA protein-group headline, the reanalysis recovered more protein groups under the original studies' $\geq 2$-unique-peptide criterion, applied to the deposited count and to the \quantmsdiann{} count under its own global protein-group rule: NCI-60 6{,}812 versus 4{,}284 ($+59$\%) and ProCan 8{,}993 versus 6{,}698 ($+34$\%; Fig.~\ref{fig:reanalysis}a). This $\geq 2$-peptide gain exceeds the change in the total protein-group count: current \texttt{DIA-NN} releases identify more peptides per protein, so proteins that the original analyses supported with a single peptide now clear the stricter two-peptide threshold, expanding the high-confidence protein set rather than only the marginal one. The Sun PCT-SWATH panel, whose deposited OpenSWATH analysis reports proteotypic proteins rather than a $\geq 2$-peptide count, shows a comparable improvement at the standard 1\% protein-FDR rule: 8{,}304 \quantmsdiann{} protein groups (\texttt{Lib.PG.Q.Value}~$\leq 0.01$) versus 6{,}183 in the deposited analysis ($+34$\%). The size of the gain reflects the age and engine of the deposited analysis, largest for the oldest, OpenSWATH-era NCI-60 ($+59$\%); the return is therefore greatest exactly where deposits are most numerous: the large back-catalogue of DIA data processed with older or non-\texttt{DIA-NN} engines. Each original count is recomputed from public data (NCI-60 and Sun from the deposited PRIDE OpenSWATH analyses, ProCan from the authors' figshare peptide-count matrix), so every bar is reproducible rather than a transcribed paper headline (Experimental Procedures). (PXD041421 and PXD017199 carry no published DIA headline and contribute reanalysis-only counts.)
Across the five cohorts the reanalysis spans 12{,}914 unique UniProt accessions, of which \textasciitilde 44\% form a pan-cohort core proteome detected in all five, resolved across 29 tissue categories (Supplementary Fig.~S6). Because every count comes from a single SDRF-parameterised invocation, the resulting protein matrix is queryable directly from the published \qpx{} archive and can be filtered or re-grouped by Cellosaurus, NCIt, or Disease Ontology terms without manual harmonisation, turning independently deposited cohorts into a single reusable cross-study resource.
The same pattern extends beyond bulk cell lines to another acquisition mode: a spatial Deep Visual Proteomics cohort of MYCN-amplified neuroblastoma yielded more identifications than its original \texttt{DIA-NN}~1.8.1 analysis (PXD064049; precursors 21{,}535 versus 17{,}958, $+20$\%; protein groups 3{,}301 versus 2{,}947, $+12$\%, both counted from the per-precursor reports under the same 1\% \texttt{Lib} q-value rather than the deposited count matrices, which are written at different run-specific q-values per \texttt{DIA-NN} version). Both sides count the identical 12 quantified runs. Notably, the deposited \texttt{DIA-NN}~1.8.1 analysis additionally loaded runs beyond these 12, which would tend to raise its globally identifiable protein count; even with that advantage, the 2.5.1 reanalysis identified more. Residual protein-group inference inflation is also higher in the deposited 1.8.1 analysis (3.2\% multi-accession groups) than in the reanalysis (1.3\%), so the protein-group gain is a conservative estimate. Per-cohort peptide-per-protein distributions and accession overlap for the reanalysed cohorts are provided in Supplementary Note~6 (Figs.~S4--S6).
%% [end results]

\section{Discussion}
\label{sec:discussion}
%% [begin sections/discussion.tex]
Public proteomics archives such as PRIDE and iProX are increasingly functioning as discovery resources, not merely as static collections of datasets deposited to support individual publications. Reuse of repository data can increase statistical power, connect independent cohorts, and test biological signals across laboratories~\cite{perezriverol2022reuse,hewapathirana2026datareuse}; workflow-based reanalysis has already shown how public deposits can be converted into reusable quantitative resources~\cite{dai2024quantms}, and resources such as ProteomicsDB show the value of aggregating protein-expression evidence across studies~\cite{lautenbacher2022proteomicsdb}. In particular, DIA datasets are becoming an increasingly valuable resource for biological discovery as their volume in public archives grows~\cite{perezriverol2025pride}, and \texttt{DIA-NN} has become one of the central engines for their analysis~\cite{demichev2020}. \quantmsdiann{} addresses this need with a scalable, SDRF-driven Nextflow workflow that parametrises every run from its own metadata and runs any supported \texttt{DIA-NN} release from a container-pinned profile.

Because the per-file work is parallelised across cloud or HPC nodes, this design makes archive-scale reanalysis practical: a 2{,}300-run single-cell cohort was reanalysed in 2.2~hours on 300 nodes, and the 5{,}798-run ProCan cohort finished in 6.2~h (Fig.~\ref{fig:validation}a,b). Across ProteoBench~\cite{devreese2025proteobench}, distributed \quantmsdiann{} reproduces the quantification of the corresponding single-machine \texttt{DIA-NN} submission (Fig.~\ref{fig:validation}c), so distributing the analysis across a cluster is an infrastructure choice rather than a change in the result. Reproducibility also holds across releases: because every \texttt{DIA-NN} version runs as a container-pinned profile driven by the same SDRF, a reanalysis can be reproduced exactly or repeated on a later release without reinstallation or reconfiguration. To our knowledge this is the first pipeline-level treatment of reproducible, version-pinned \texttt{DIA-NN} reanalysis, which grows in importance as the engine evolves and archived data are revisited.

Beyond reproducing the original analysis, reanalysis with the current \texttt{DIA-NN} builds used here adds depth, an effect already reported when public data were reprocessed with quantms~\cite{dai2024quantms}. The gain is clearest in low-input single-cell data, where match-between-runs matters most: upgrading the engine raised single-cell precursor and protein-group identifications by up to \textasciitilde 17\%. Bulk cell-line deposits processed with older or non-\texttt{DIA-NN} engines show the same effect at larger magnitude, with reanalysis recovering 34\% to 59\% more protein groups per cohort (for example, NCI-60 $+59$\% and ProCan $+34$\%) and the largest gains on the oldest deposits. The largest recoverable gains therefore lie in the older DIA back-catalogue, where many datasets were processed with earlier software generations.

Depth and speed are only useful if the results can be interpreted and trusted, and biological interpretation requires knowing what each sample is, how it was acquired, and how those attributes map onto the quantification matrix; the scarcity of structured sample metadata remains a major barrier to proteomics reuse~\cite{perezriverol2020metadata,perezriverol2022reuse,deutsch2026proteomexchange}. \quantmsdiann{} treats the SDRF~\cite{dai2021sdrf} as both the analytical control layer and the annotation source: every per-run \texttt{DIA-NN} setting is derived from it, and the same file annotates the resulting \qpx{} archive, so identifications, quantifications, and provenance travel together in one queryable artefact. This is what allowed the five-cohort cell-line reanalysis to be regrouped by Cellosaurus, NCIt, or Disease Ontology terms directly from the published output, with no manual harmonisation. Because each run is parametrised from its own metadata, the same design spans acquisition modes without modification, from single-cell to plexDIA, spatial Deep Visual Proteomics, and phosphoproteomics~\cite{brunner2022,derks2023,mund2022,makhmut2023}, and other \texttt{DIA-NN} applications such as immunopeptidomics~\cite{scheid2025mhcquant2} follow the same path. This combination of current \texttt{DIA-NN} releases, SDRF-first configuration, scalable execution, and standardised \qpx{} output extends earlier DIA workflows~\cite{bichmann2021diaproteomics}.

\quantmsdiann{} therefore provides actively maintained infrastructure for DIA meta-analysis and archive-scale proteomics reuse. Its value is amplified by the broader quantms ecosystem (\url{https://quantms.org}) rather than ending at a single workflow: \quantmsdiann{} produces harmonised \qpx{} and MSstats-ready outputs, \texttt{pmultiqc} provides metadata-aware quality control~\cite{yue2026pmultiqc}, downstream single-cell ecosystems support batch correction, cell-type assignment, and differential analysis~\cite{vanderaa2023,virshup2023scverse}, and the sibling quantms workflow extends the same SDRF/nf-core model to large-scale quantitative proteomics beyond DIA, including DDA and TMT-style analyses~\cite{dai2024quantms}. Applied across the growing public archive, this ecosystem can integrate independently deposited cohorts into reusable resources and move DIA data towards protein-expression knowledge bases~\cite{lautenbacher2022proteomicsdb} built from the community's accumulated measurements~\cite{perezriverol2025pride}. \quantmsdiann{} is open-source under the MIT license and developed within the nf-core and bigbio proteomics ecosystem, allowing it to track new \texttt{DIA-NN} releases, acquisition modes, and metadata requirements as they appear.

%% [end discussion]

%% [begin sections/declarations.tex]
\section*{Data Availability}
All datasets reanalysed in this study are publicly available through ProteoMeXchange.
Cell-line bulk DIA reanalyses: PRIDE PXD003539 (Guo 2019, NCI-60), PXD004701 (Sun 2023), PXD030304 (ProCan-DepMapSanger), PXD017199 (Tognetti 2021), PXD041421 (Wang 2023). Spatial Deep Visual Proteomics reanalysis: PRIDE PXD064049 (MYCN-amplified neuroblastoma, diaPASEF).
Scalability experiment: PRIDE PXD071075 (Wu 2025). ProteoBench benchmarks: PXD049412 (single-cell, Module 9), PXD062685 (diaPASEF, Module 5), PXD070049 (ZenoTOF, Module 10), and Module 7 (Astral) datasets. The full list of PRIDE accessions used in the pan-cohort of DIA cell-line reanalyses is provided in Supplementary Note~2 (Table~S1). \quantmsdiann{} is available at \url{https://github.com/bigbio/quantmsdiann} under the MIT license; documentation is hosted at \url{https://quantmsdiann.quantms.org/}.
The pipeline is implemented in Nextflow DSL2 ($\geq$25.10.4) and packages each step as a versioned Docker, Singularity, or Apptainer container image. \texttt{DIA-NN}~1.8.1 is distributed publicly via BioContainers; because the \texttt{DIA-NN} license restricts redistribution of later releases, containers for \texttt{DIA-NN}~1.9 and above are built locally from the recipes in \url{https://github.com/bigbio/quantms-containers} (see Experimental Procedures). Pipeline version v2.2.0 (release codename ``Chongqing'', 2026-06-16) was used for the analyses reported in this paper. Analysis scripts that generated the figures in this manuscript are available at \url{https://github.com/ypriverol/quantmsdiann-paper} (this repository).

\section*{Supplemental Data}
This article contains supplemental data.
The Supplementary Notes provide the \texttt{DIA-NN} container-build recipes (Supplementary Note~1), the full PRIDE accession list (Supplementary Note~2, Table~S1), the filtering and protein-counting rules (Supplementary Note~3), per-process runtime and resource envelopes (Supplementary Note~4, Figs.~S1--S2), the ProteoBench per-version identification breakdown (Supplementary Note~5, Fig.~S3), per-cohort peptide-per-protein and accession-overlap figures (Supplementary Note~6, Figs.~S4--S6), and a table of software, formats, and public resources with their repositories (Supplementary Note~7, Table~S2).

\section*{Conflict of Interest}
V.D.\ holds shares of Aptila Biotech. The other authors declare no competing interests.

\section*{Acknowledgments}
M.B.\ is supported by the Science and Technology Innovation Key R\&D Program of Chongqing (CSTB2023TIAD-STX0002). A.L.\ acknowledges support from MCIN (Recovery, Transformation and Resilience Plan); the Basque Government ``Biotechnology Complementary Plan Applied to Health'' funded by the European Union NextGeneration EU (PRTR-C17.I1; PRTR-C17.I01.P01.S13) (AAAA\_ACG\_AY\_2539/22\_05); and an EMBO Scientific Exchange Grant (number 10015/2023). Y.P-R.\ acknowledges support from the Wellcome Trust [223745/Z/21/Z], BBSRC [BB/V018779/1, BB/X001911/1], and EPSRC [EP/Y035984/1]. Y.P-R., N.S., O.K.\ and J.A.V.\ acknowledge support from a UKRI single-cell proteomics (SCP) grant [UKRI701]. Y.P-R.\ and J.A.V.\ acknowledge EMBL core funding. V.D.\ was funded by the German Ministry of Education and Research (BMBF), as part of the National Research Node ``Mass spectrometry in Systems Medicine'' (MSCoreSys), under grant agreement 161L0221. We thank the EMBL-EBI HPC team for cluster support, and the ProteoBench community for providing a reproducible benchmark of DIA search engines.

\section*{Use of Generative AI}
During the preparation of this manuscript, the authors used generative AI tools (Anthropic Claude, and OpenAI Codex and ChatGPT) to assist with software code generation and manuscript editing. All content was reviewed and edited by the authors, who take full responsibility for it.

\section*{Author Contributions}
QXY, YS, CD and Y.P-R.\ developed the pipeline. HW helped build the containers and contributed to the workflow implementation and manuscript writing. AL annotated the SDRF metadata and helped run the workflow.
TS, FC, NS and JN assisted with the DIA-NN reanalyses and manuscript writing; TS additionally contributed to the workflow design and to the design of \qpx{} and the SDRF format, and JN tested the tool and identified multiple issues. OK provided feedback on the plexDIA datasets. VD provided DIA-NN expertise and guidance on cross-version reproducibility. JAV provided proteomics data-resource infrastructure and funding, supervised the EMBL-EBI contributors, and contributed to manuscript review and editing. MB contributed to manuscript writing and supervised the students; MB and Y.P-R.\ conceived and supervised the project. QXY and Y.P-R.\ wrote the original draft, with input from all authors. All authors read and approved the final manuscript.

\section*{Abbreviations}
\begin{description}[leftmargin=5.5em,style=nextline,itemsep=0pt,parsep=0pt]
\item[CLI] command-line interface
\item[DDA] data-dependent acquisition
\item[DIA] data-independent acquisition
\item[DIA-NN] data-independent acquisition neural networks (DIA search and quantification software)
\item[diaPASEF] data-independent acquisition with parallel accumulation--serial fragmentation
\item[DVP] Deep Visual Proteomics
\item[EFO] Experimental Factor Ontology
\item[FDR] false discovery rate
\item[HPC] high-performance computing
\item[HUPO-PSI] Human Proteome Organization Proteomics Standards Initiative
\item[HYE] human, yeast, and \textit{E.\ coli} (ProteoBench three-proteome mixture)
\item[IQR] interquartile range
\item[LSF] Load Sharing Facility (HPC scheduler)
\item[MS] mass spectrometry
\item[MS1/MS2] first-/second-stage mass spectrometry (precursor/fragment scans)
\item[PASEF] parallel accumulation--serial fragmentation
\item[PSI-MS] HUPO-PSI mass spectrometry controlled vocabulary
\item[PSM] peptide-spectrum match
\item[QPX] Quantitative Proteomics eXchange
\item[SDRF] Sample and Data Relationship Format
\end{description}
%% [end declarations]

\bibliographystyle{elsarticle-num}
\bibliography{references}

\end{document}
